\section{Related Work}
\label{sec:related_work}

\paragraph{Reasoning supervision and trace selection.}
Chain-of-thought reasoning exposes intermediate steps that can serve as supervision for smaller models~\citep{wei2022chain,hsieh2023distilling}.
However, stronger teachers and longer traces do not consistently improve small-student performance; mixing reasoning of different complexity can be more effective~\citep{li2025small}.
This motivates comparisons among random-correct, shortest-correct, and mixed-length supervision from a common candidate pool.
Our within-problem design fixes teacher identity, problem content, and final-answer correctness, while separately accounting for supervision volume.
It investigates the usefulness of generated sequences through SFT, without teacher-logit matching.

\paragraph{Compression before distillation.}
LiteCoT uses difficulty-aware prompting to rewrite teacher solutions into concise traces~\citep{wu2025concise}.
Compress-Distill studies post-hoc compression at multiple student scales and reports lower training and inference costs alongside an accuracy trade-off~\citep{griot2026compress}.
These methods establish textual compression as a direct comparator for constructing student supervision.
Our framework modifies activations while the teacher generates a solution, but this implementation difference alone does not establish better teaching data.
Comparisons must account for reference-trace generation, rewriting calls, usable-answer coverage, and student training budgets.
An ordinary shortening prompt is also distinct from a full difficulty-aware rewriting pipeline and should be evaluated under its own name.

\paragraph{Activation-based reasoning control.}
ASC learns a steering vector from verbose--concise pairs to shorten chain-of-thought generation~\citep{azizi2026activation}.
Its published formulation uses a length-normalized contrastive energy objective and a KL trust-region constraint, making it distinct from a simple difference-of-means baseline.
Other work extracts SAE features to steer reasoning and also studies an SAE-free alternative~\citep{li2025feature}, or identifies strategy-associated features through keyword screening followed by intervention-based ranking~\citep{fang2026controllable}.
Our focus is how naturally occurring length contrasts inform the generation of student training data.
Dense directions, concise prompting, and norm-matched random feature sets test whether an SAE dictionary provides incremental value for this objective.
Because our selected features have high activation near answer markers, an answer-format direction is an additional mechanistic control specific to the present evidence.

\paragraph{Sparse features, interpretation, and transfer.}
SAEs provide sparse decompositions of model activations, with TopK architectures controlling the number of active features~\citep{cunningham2023sparse,gao2024scaling}.
Such decompositions do not guarantee that individual features correspond to reasoning operations.
\citet{ma2026reasoning} demonstrate that contrastively identified candidates can respond to lexical injections and non-reasoning counterexamples, motivating tests of semantic invariance and alternative explanations.
We therefore separate observed length associations, intervention-induced changes to model states, changes to generated content, and effects on student learning.
Steering vector distillation already demonstrates behavioral transfer through outputs from steered teachers~\citep{blank2026subliminal}; the teacher-intervention-to-student-training pipeline itself is not our novelty claim.
Our proposed baseline and SAE analyses instead test what the selected directions control and whether that control provides useful supervision beyond simpler ways of shortening responses.
